% =============================================================================
%  lec21.tex —— 第 21 讲 $\infty$-操作子（底本 PDF 第 161–168 页）
% =============================================================================
\lecture{21}{$\infty$-操作子}
\begin{center}\small\color{black!55}（底本 PDF 第 161–168 页）\end{center}

本讲介绍 $\infty$-操作子与幺半 $\infty$-范畴。

\section{对称幺半 $\infty$-范畴}

\nitem\label{L21-1.1} 回忆，一个\textbf{幺半范畴}是配备了下述结构的（普通）范畴 $\cC$：
\begin{itemize}[leftmargin=2.2em,itemsep=0.1em]
\item 一个二元函子 $-\otimes-\colon\cC\times\cC\to\cC$，称为\textbf{幺半积}；
\item 一个对象 $\mathbf{1}\in\cC$，称为\textbf{幺半单位}；
\item 一个形如 $X\otimes(Y\otimes Z)\simeq(X\otimes Y)\otimes Z$ 的可逆自然变换，称为\textbf{结合子}；
\item 形如 $\mathbf{1}\otimes X\simeq X$、$X\otimes\mathbf{1}\simeq X$ 的可逆自然变换，称为\textbf{左、右单位子}，
\end{itemize}
使得包括五边形图在内的某些图交换。

非形式地说，我们可以把幺半范畴说成是一个配备了乘法、且该乘法在连贯同伦意义下具有单位性与结合性的范畴。这里连贯数据是有限的，因为普通范畴构成一个 $2$-范畴。事实上，幺半范畴的定义只涉及 $0\le n\le4$ 的 $n$ 元运算 $\cC^{\times n}\to\cC$。

\nitem\label{L21-1.2} 同样回忆，一个\textbf{对称幺半范畴}是配备了下述结构的幺半范畴 $\cC$：
\begin{itemize}[leftmargin=2.2em,itemsep=0.1em]
\item 一个形如 $X\otimes Y\simeq Y\otimes X$ 的可逆自然变换，称为\textbf{换位自然变换}，
\end{itemize}
使得包括逆律在内的某些图交换。

\nitem\label{L21-1.3} 我们希望把上述内容推广为（对称）幺半 $\infty$-范畴的概念。此时连贯数据应当涉及所有 $n\ge0$ 的 $n$ 元运算。

我们不想把这些连贯数据写下来，而希望把它们编码为一个函子 $F\colon J\to\Cat_{\infty}$，又称 $\infty$-范畴的同伦连贯图，其中 $J$ 是一个精心设计的单纯集，使得
\begin{itemize}[leftmargin=2.2em,itemsep=0.1em]
\item 对 $n\ge0$，有一个顶点 $\langle n\rangle$，它被 $F$ 送到 $n$ 次幂 $\cC^{\times n}$；
\item 有一条边 $\langle 2\rangle\to\langle 1\rangle$，编码幺半积；
\item 有一条边 $\langle 0\rangle\to\langle 1\rangle$，编码幺半单位；
\item 连贯数据由 $J$ 中的胞腔编码。
\end{itemize}
注意，第一点实际上是结构而非性质：我们需要对每个 $1\le k\le n$ 提供投影 $p_k\colon F(\langle n\rangle)\to F(\langle 1\rangle)$，它们把 $F(\langle n\rangle)$ 实现为 $F(\langle 1\rangle)$ 的 $n$ 次幂。因此我们应当有：
\begin{itemize}[leftmargin=2.2em,itemsep=0.1em]
\item 对每个 $i\in\{1,\cdots,n\}$，有一条边 $\langle n\rangle\to\langle 1\rangle$，编码投影函子 $\mathrm{pr}_i\colon\cC^{\times n}\to\cC$。
\end{itemize}
此外，在对称幺半情形，为了编码换位自然变换，我们需要能够交换 $\cC^{\times n}$ 中的因子。因此我们应当有：
\begin{itemize}[leftmargin=2.2em,itemsep=0.1em]
\item 对 $n\ge0$，对称群 $\Sigma_n$ 应当作用在 $\langle n\rangle$ 上。
\end{itemize}

事实上，上述两条要求可以合并为：
\begin{itemize}[leftmargin=2.2em,itemsep=0.1em]
\item 对每个单射 $\varphi\colon\{1,\cdots,m\}\to\{1,\cdots,n\}$，有一条边 $\langle n\rangle\to\langle m\rangle$，编码投影函子 $\mathrm{pr}_{\varphi}\colon\cC^{\times n}\to\cC^{\times m}$。
\end{itemize}
$J$ 中的这些边将称为\textbf{惰性态射}。相反，编码乘积与单位等数据的边称为\textbf{活动态射}。

我们的直觉是：
\begin{itemize}[leftmargin=2.2em,itemsep=0.1em]
\item 在对称幺半情形，存在唯一的活动态射 $\langle n\rangle\to\langle 1\rangle$。
\item 在幺半情形，$\Sigma_n$ 应当自由且传递地作用在活动态射 $\langle n\rangle\to\langle 1\rangle$ 上。
\end{itemize}

\begin{exercise}\label{L21-1.4}
证明：单射 $\varphi\colon\{1,\cdots,m\}\to\{1,\cdots,n\}$ 与满足下述条件的带标记集映射
\[
\alpha\colon\{*,1,2,\cdots,n\}\to\{*,1,2,\cdots,m\}
\]
自然双射：对任意 $1\le j\le m$，$\alpha^{-1}(\{j\})$ 是单点集。
\end{exercise}

\nitem\label{L21-1.5} 上述讨论引出 Segal 的下述定义。

\begin{definition}\label{L21-1.6}
设 $\mathrm{Fin}_{*}$ 是带标记有限集构成的普通范畴的极小模型。对每个 $n\ge0$，令 $\langle n\rangle\in\mathrm{Fin}_{*}$ 为对象
\[
\langle n\rangle\defeq\{*\}\sqcup\langle n\rangle^{\circ}\defeq\{*,1,2,\cdots,n\}.
\]
设 $\alpha\colon\langle n\rangle\to\langle m\rangle$ 是 $\mathrm{Fin}_{*}$ 中的态射，我们说
\begin{enumerate}[label=(\arabic*),leftmargin=2.4em,itemsep=0.1em]
\item 态射 $\alpha$ 是\textbf{惰性的}，若对任意非标记元素 $j\in\langle m\rangle^{\circ}$，$\alpha^{-1}(\{j\})$ 是单点集。
\item 态射 $\alpha$ 是\textbf{活动的}，若 $\alpha^{-1}(\{*\})=\{*\}$。
\end{enumerate}
我们还记
\[
\mathrm{Comm}^{\otimes}\defeq\mathrm{Fin}_{*}.
\]
\end{definition}

\begin{exercise}\label{L21-1.7}
$\mathrm{Comm}^{\otimes}$ 中的态射是同构，当且仅当它既是惰性的又是活动的。
\end{exercise}

\begin{exercise}\label{L21-1.8}
任意态射 $\langle n\rangle\to\langle m\rangle$ 都可以写成
\[
\langle n\rangle\xrightarrow{\ \beta\ }\langle l\rangle\xrightarrow{\ \gamma\ }\langle m\rangle,
\]
其中 $\beta$ 是惰性的、$\gamma$ 是活动的。此外，这样的表达式在唯一同构的意义下唯一。
\end{exercise}

\begin{definition}\label{L21-1.9}
一个\textbf{对称幺半 $\infty$-范畴}是函子
\[
F\colon\mathrm{Comm}^{\otimes}\to\Cat_{\infty},
\]
它把惰性态射 $\rho^{i}\colon\langle n\rangle\to\langle 1\rangle$（$i\in\langle n\rangle^{\circ}$）送到函子
\[
F(\rho^{i})\colon F(\langle n\rangle)\to F(\langle 1\rangle),
\]
这些函子把 $F(\langle n\rangle)$ 实现为 $F(\langle 1\rangle)^{\times n}$。我们称 $F(\langle 1\rangle)$ 为它的\textbf{底层} $\infty$-范畴。
\end{definition}

\begin{exercise}\label{L21-1.10}
若 $F(\langle 1\rangle)$ 是普通的，则 $F$ 决定一个对称幺半范畴。
\end{exercise}

\nitem\label{L21-1.11} 非对称情形可以用 $\mathrm{Comm}^{\otimes}$ 的下述变体来编码。

\begin{definition}\label{L21-1.12}
设 $\mathrm{Assoc}^{\otimes}$ 是按如下方式定义的普通范畴：
\begin{itemize}[leftmargin=2.2em,itemsep=0.1em]
\item $\mathrm{Assoc}^{\otimes}$ 中的一个对象是带标记有限集 $\langle n\rangle$，其中 $n\ge0$。
\item 从 $\langle n\rangle$ 到 $\langle m\rangle$ 的一个态射是一个对 $(\alpha,\preceq_j)$，其中
\begin{itemize}[leftmargin=1.9em,itemsep=0.05em]
\item $\alpha\colon\langle n\rangle\to\langle m\rangle$ 是带标记集之间的映射；
\item 对每个 $j\in\langle m\rangle^{\circ}$，$\preceq_j$ 是 $\alpha^{-1}(j)$ 上的线性序。
\end{itemize}
\end{itemize}
对于态射 $(\alpha,\preceq_j)$，若 $\alpha$ 是惰性的（相应地，活动的），则称它是惰性的（相应地，活动的）。
\end{definition}

\begin{definition}\label{L21-1.13}
一个\textbf{幺半 $\infty$-范畴}是函子
\[
F\colon\mathrm{Assoc}^{\otimes}\to\Cat_{\infty},
\]
它把惰性态射 $\rho^{i}\colon\langle n\rangle\to\langle 1\rangle$（$i\in\langle n\rangle^{\circ}$）送到函子
\[
F(\rho^{i})\colon F(\langle n\rangle)\to F(\langle 1\rangle),
\]
这些函子把 $F(\langle n\rangle)$ 实现为 $F(\langle 1\rangle)^{\times n}$。我们称 $F(\langle 1\rangle)$ 为它的\textbf{底层} $\infty$-范畴。
\end{definition}

\begin{exercise}\label{L21-1.14}
证明：练习~\ref{L21-1.7}--\ref{L21-1.10} 在把 $\mathrm{Fin}_{*}$ 换成 $\mathrm{Assoc}^{\otimes}$ 后仍然成立。
\end{exercise}

\begin{exercise}\label{L21-1.15}
证明：对 $\mathrm{Comm}^{\otimes}$ 与 $\mathrm{Assoc}^{\otimes}$ 两者，$\langle n\rangle$ 上的自同构群都可以等同于 $\Sigma_n$。
\end{exercise}

\nitem\label{L21-1.16} 存在一个函子 $\mathrm{Assoc}^{\otimes}\to\mathrm{Comm}^{\otimes}$，它通过限制把对称幺半 $\infty$-范畴送到幺半 $\infty$-范畴。

\begin{warning}\label{L21-1.17}
尽管 $\mathrm{Assoc}^{\otimes}\to\mathrm{Comm}^{\otimes}$ 的行为像商函子（它在对象与态射上都满射），但不应把对称幺半 $\infty$-范畴想成是满足某些性质的幺半 $\infty$-范畴。更确切地说，即便对于取值于 $(2,1)$-范畴 $D$（例如普通范畴构成的 $(2,1)$-范畴）的函子 $\mathrm{Assoc}^{\otimes}\to D$，穿过 $\mathrm{Comm}^{\otimes}$ 的弱分解也是结构而非性质。
\end{warning}

\begin{remark}\label{L21-1.18}
事实上，我们可以用更高效的方式编码非对称幺半 $\infty$-范畴：把 $\mathrm{Assoc}^{\otimes}$ 的角色换成 $\Delta^{\op}$。这称为幺半群的单纯模型。粗略地说，这相当于从函子 $\mathrm{Assoc}^{\otimes}\to\Cat_{\infty}$ 所编码的数据中消去 $\Sigma_n$-作用 $\Hom(\langle n\rangle,\langle 1\rangle)$。

与 $\mathrm{Assoc}^{\otimes}$ 相比，这个单纯模型有几个优点。例如，可以用截断 $(\Delta_{\le n})^{\op}$ 定义 $A_n$-代数。然而，由于 $\Sigma_n$-作用被隐藏了起来，很难把单纯模型编码的幺半结构与对称幺半结构直接比较。

详见 HA.4.1 与 4.2.2。
\end{remark}

\section{$\infty$-操作子}

\begin{exercise}\label{L21-2.1}
证明：对 $\mathrm{Comm}^{\otimes}$ 与 $\mathrm{Assoc}^{\otimes}$ 两者，下述数据等价：
\begin{itemize}[leftmargin=2.2em,itemsep=0.1em]
\item 一个活动态射 $\alpha\colon\langle n\rangle\to\langle m\rangle$；
\item 一个分解
\[
\langle n\rangle^{\circ}\simeq\bigsqcup_{j\in\langle m\rangle^{\circ}}\langle n_j\rangle^{\circ},
\]
以及一族以 $j\in\langle m\rangle^{\circ}$ 为指标的活动态射 $\alpha_j\colon\langle n_j\rangle\to\langle 1\rangle$。
\end{itemize}
提示：对每个 $\alpha$，考虑练习~\ref{L21-1.8} 中 $\rho_j\circ\alpha$ 的分解。
\end{exercise}

\nitem\label{L21-2.2} 事实上，范畴 $\mathrm{Comm}^{\otimes}$ 与 $\mathrm{Assoc}^{\otimes}$ 可以由下述数据还原：
\begin{enumerate}[label=(\roman*),leftmargin=2.4em,itemsep=0.1em]
\item 对每个 $n\ge0$，活动态射 $\langle n\rangle\to\langle 1\rangle$ 的集合 $\mathrm{O}(n)$；
\item $\Sigma_n$ 在 $\mathrm{O}(n)$ 上的作用；
\item 对 $m\ge0$ 以及以 $j\in\langle m\rangle^{\circ}$ 为指标的 $n_j$，一个映射
\[
\mathrm{O}(m)\times\Bigl(\prod_j\mathrm{O}(n_j)\Bigr)\to\mathrm{O}(n_1+\cdots+n_m),
\]
它由上述练习给出。
\end{enumerate}
在经典文献中，一族满足某些相容性的这类数据称为一个（对称）\textbf{操作子}。这里 $\mathrm{O}(n)$ 是 $n$ 元运算的集合，而 (iii) 中的映射称为\textbf{复合映射}。

\nitem\label{L21-2.3} 人们也可以考虑带有充实结构的操作子。例如，通过在以普通对称幺半范畴 $\Set_{\Delta}$、$\Top$、$\Mod_k^{\heartsuit}$、$\Ch(\Ab)$ 等上的充实，我们有单纯操作子、拓扑操作子、$k$-线性操作子、dg-操作子等概念。

我们希望有一个在 $\infty$-范畴 $\Spc$ 上弱充实的操作子概念。换言之，每个 $\mathrm{O}(n)$ 都应当是一个空间，而复合映射在连贯同伦的意义下结合。

遵循本课程的一般哲学，更好的做法是把 $\infty$-操作子定义为满足某些条件的 $\infty$-范畴（例如 $\mathrm{Comm}^{\otimes}$ 与 $\mathrm{Assoc}^{\otimes}$），而不是罗列所有这些连贯数据。

\nitem\label{L21-2.4} 在给出 $\infty$-操作子的定义之前，让我们再做一次推广。

设 $A$ 是对称幺半 $\infty$-范畴，$M$ 是 $\infty$-范畴。我们想通过函子 $F\colon\mathrm{CMod}^{\otimes}\to\Cat_{\infty}$ 定义 $A$ 在 $M$ 上的作用。指标范畴 $\mathrm{CMod}^{\otimes}$ 应当以 $\mathrm{Comm}^{\otimes}$ 为全子范畴，其余部分则应当编码 $M$ 上的 $A$-模结构。因此至少应当有：
\begin{itemize}[leftmargin=2.2em,itemsep=0.1em]
\item 来自 $\langle n\rangle\in\mathrm{Assoc}^{\otimes}$ 的对象 $(\mathbf{a},\cdots,\mathbf{a})$，它们应当被 $F$ 送到 $A^{\times n}$；
\item 对象 $(\mathbf{a},\cdots,\mathbf{a},\mathbf{m})$，它们应当被 $F$ 送到 $A^{\times n}\times M$；
\item 唯一的活动态射 $(\mathbf{a},\cdots,\mathbf{a},\mathbf{m})\to\mathbf{m}$，编码作用函子 $A^{\times n}\times M\to M$。
\end{itemize}
注意我们还应当有形如 $(\mathbf{m},\mathbf{a},\cdots,\mathbf{a})$ 的对象，更一般地形如 $(\mathbf{a},\cdots,\mathbf{a},\mathbf{m},\mathbf{a},\cdots,\mathbf{a})$ 的对象。出于技术原因，更方便的做法是甚至允许含多个 $\mathbf{m}$ 项的顺序列。当然，我们不会允许任何从 $(\mathbf{m},\mathbf{m})$ 到 $\mathbf{m}$ 的活动态射。

\begin{definition}\label{L21-2.5}
设 $\mathrm{CMod}^{\otimes}$ 是按如下方式定义的范畴：
\begin{itemize}[leftmargin=2.2em,itemsep=0.1em]
\item 一个对象是一个对 $(\langle n\rangle,c)$，其中 $\langle n\rangle\in\mathrm{Fin}_{*}$，$c\colon\langle n\rangle^{\circ}\to\{\mathbf{a},\mathbf{m}\}$ 是映射。
\item 从 $(\langle n\rangle,c)$ 到 $(\langle m\rangle,d)$ 的一个态射是 $\mathrm{Comm}^{\otimes}$ 中的态射 $\alpha\colon\langle n\rangle\to\langle m\rangle$，使得
\begin{itemize}[leftmargin=1.9em,itemsep=0.05em]
\item 若 $c(i)=\mathbf{m}$，则 $\alpha(i)=*$ 或 $d(\alpha(i))=\mathbf{m}$；
\item 若 $d(j)=\mathbf{m}$，则存在唯一的 $i\in\alpha^{-1}(\{j\})$ 使 $c(i)=\mathbf{m}$。
\end{itemize}
\end{itemize}
对于这样的态射，若底层的 $\mathrm{Comm}^{\otimes}$ 中的态射是惰性的（相应地，活动的），则称它是惰性的（相应地，活动的）。
\end{definition}

\nitem\label{L21-2.6} 类似地，我们可以定义编码幺半群及其左模的范畴 $\mathrm{LMod}^{\otimes}$。

\begin{definition}\label{L21-2.7}
设 $\mathrm{LMod}^{\otimes}$ 是按如下方式定义的范畴：
\begin{itemize}[leftmargin=2.2em,itemsep=0.1em]
\item 一个对象是一个对 $(\langle n\rangle,c)$，其中 $\langle n\rangle\in\mathrm{Fin}_{*}$，$c\colon\langle n\rangle^{\circ}\to\{\mathbf{a},\mathbf{m}\}$ 是映射。
\item 从 $(\langle n\rangle,c)$ 到 $(\langle m\rangle,d)$ 的一个态射是 $\mathrm{Assoc}^{\otimes}$ 中的态射 $(\alpha,\preceq_j)\colon\langle n\rangle\to\langle m\rangle$，使得
\begin{itemize}[leftmargin=1.9em,itemsep=0.05em]
\item 若 $c(i)=\mathbf{m}$，则 $\alpha(i)=*$ 或 $d(\alpha(i))=\mathbf{m}$；
\item 若 $d(j)=\mathbf{m}$，则存在唯一的 $i\in\alpha^{-1}(\{j\})$ 使 $c(i)=\mathbf{m}$，且 $i$ 也是关于序 $\preceq_j$ 的最大元。
\end{itemize}
\end{itemize}
对于这样的态射，若底层的 $\mathrm{Assoc}^{\otimes}$ 中的态射是惰性的（相应地，活动的），则称它是惰性的（相应地，活动的）。

对偶地，我们把“最大”换成“最小”，定义 $\mathrm{RMod}^{\otimes}$。
\end{definition}

\nitem\label{L21-2.8} 诸如 $\mathrm{CMod}^{\otimes}$、$\mathrm{LMod}^{\otimes}$ 或 $\mathrm{RMod}^{\otimes}$ 这样的范畴是\textbf{有色操作子}，而符号 $\mathbf{a}$、$\mathbf{m}$ 是颜色。

\begin{exercise}\label{L21-2.9}
构造一个带三个颜色 $(\mathbf{a}_l,\mathbf{m},\mathbf{a}_r)$ 的普通有色操作子 $\mathrm{BMod}^{\otimes}$，它编码两个幺半群以及它们的一个双模。
\end{exercise}

\begin{exercise}\label{L21-2.10}
设 $\OO^{\otimes}$ 是 $\mathrm{Comm}^{\otimes}$、$\mathrm{Assoc}^{\otimes}$、$\mathrm{CMod}^{\otimes}$、$\mathrm{LMod}^{\otimes}$、$\mathrm{RMod}^{\otimes}$ 或 $\mathrm{BMod}^{\otimes}$ 之一。考虑自然函子 $p\colon\OO^{\otimes}\to\mathrm{Fin}_{*}$。证明：
\begin{enumerate}[label=(\arabic*),leftmargin=2.4em,itemsep=0.1em]
\item 对 $\mathrm{Fin}_{*}$ 中的任意惰性态射，它都有足够多的 $p$-余笛卡尔提升。此外，$\OO^{\otimes}$ 中的一个态射是惰性的，当且仅当它是 $\mathrm{Fin}_{*}$ 中某个惰性态射上方的 $p$-余笛卡尔态射。
\item 设 $C\in\OO^{\otimes}_{\langle n\rangle}$ 是位于 $\langle n\rangle$ 之上的对象，$C\to C_i$ 是位于 $\rho^{i}\colon\langle n\rangle\to\langle 1\rangle$ 之上的 $p$-余笛卡尔态射，$i\in\langle n\rangle^{\circ}$。则这些态射把 $C$ 实现为一个 $p$-极限\footnote{关于相对极限的一般讨论，见 HTT.4.3.2。}。换言之，对任意测试对象 $D\in\OO^{\otimes}$，下述方块是笛卡尔方块：
\[
\begin{tikzcd}[column sep=4.8em,row sep=3.2em]
\mathrm{Maps}_{\OO^{\otimes}}(D,C)\arrow[r]\arrow[d] & \prod_i\mathrm{Maps}_{\OO^{\otimes}}(D,C_i)\arrow[d] \\
\mathrm{Maps}_{\mathrm{Fin}_{*}}(p(D),\langle n\rangle)\arrow[r] & \prod_i\mathrm{Maps}_{\mathrm{Fin}_{*}}(p(D),\langle 1\rangle).
\end{tikzcd}
\]
\item 协变传输函子 $\rho^{i}_{\dagger}$ 把 $\OO^{\otimes}_{\langle n\rangle}$ 实现为 $(\OO^{\otimes}_{\langle 1\rangle})^{\times n}$。
\end{enumerate}
\end{exercise}

\begin{definition}\label{L21-2.11}
一个\textbf{$\infty$-操作子}\footnote{更好的术语也许是“有色 $\infty$-操作子”。}是 $\mathrm{Fin}_{*}$ 上的 $\infty$-范畴 $\OO^{\otimes}$，它满足上述练习中的结论。\footnote{更确切地说，我们应当说：函子 $p\colon\OO^{\otimes}\to\mathrm{Fin}_{*}$ 把 $\OO^{\otimes}$ 实现为一个本质 $\infty$-操作子，若 $p$ 作为内纤维化的任一/所有实现 $\OO^{\otimes}\to\mathrm{N}(\mathrm{Fin}_{*})$ 都满足上述条件。} 我们称
\[
\OO\defeq\OO^{\otimes}_{\langle 1\rangle}
\]
为 $\OO^{\otimes}$ 中颜色构成的 $\infty$-范畴。
\end{definition}

\begin{definition}\label{L21-2.12}
设 $p\colon\OO^{\otimes}\to\mathrm{Fin}_{*}$ 是 $\infty$-操作子，$f\colon C\to C'$ 是 $\OO^{\otimes}$ 中的态射，我们说
\begin{itemize}[leftmargin=2.2em,itemsep=0.1em]
\item 态射 $f$ 是\textbf{惰性的}，若它是 $\mathrm{Fin}_{*}$ 中某个惰性态射上方的 $p$-余笛卡尔态射。
\item 态射 $f$ 是\textbf{活动的}，若 $p(f)$ 是活动的。
\end{itemize}
\end{definition}

\nitem\label{L21-2.13} 我们常把对象 $C\in\OO^{\otimes}_{\langle n\rangle}$ 记作 $(C_1,\cdots,C_n)$，\footnote{Lurie 的书使用 $C_1\oplus\cdots\oplus C_n$，这可能引起混淆，因为 $C$ 既不是诸 $C_i$ 的（绝对）积，也不是它们的余积。} 其中 $C_i\in\OO$，并把惰性态射 $C\to C_i$ 视为隐含的。

\begin{exercise}\label{L21-2.14}
$\OO^{\otimes}$ 中的一个态射是同构，当且仅当它既是惰性的又是活动的。
\end{exercise}

\begin{exercise}\label{L21-2.15}
证明：$\OO^{\otimes}$ 中的任意态射都可以本质唯一地写成先一个惰性态射、后一个活动态射的复合。
\end{exercise}

\nitem\label{L21-2.16} 设 $\OO^{\otimes}$ 是 $\infty$-操作子。我们可以把一个 $\OO$-幺半 $\infty$-范畴定义为满足某些条件的函子 $\OO^{\otimes}\to\Cat_{\infty}$。在实践中，更好的做法是通过 $\OO^{\otimes}$ 上相应的余笛卡尔纤维化来描述这些条件。

\begin{definition}\label{L21-2.17}
设 $\OO^{\otimes}$ 是 $\infty$-操作子。一个 \textbf{$\OO$-幺半 $\infty$-范畴}是一个（本质）余笛卡尔纤维化 $\cC^{\otimes}\to\OO^{\otimes}$，使得复合 $\cC^{\otimes}\to\OO^{\otimes}\to\mathrm{Fin}_{*}$ 把 $\cC^{\otimes}$ 实现为 $\infty$-操作子。我们称 $\cC\defeq\cC^{\otimes}_{\langle 1\rangle}$ 为 $\cC^{\otimes}$ 的\textbf{底层} $\infty$-范畴。
\end{definition}

\begin{definition}\label{L21-2.18}
设 $\cC^{\otimes}$ 与 $\cC^{\prime\otimes}$ 是 $\OO$-幺半 $\infty$-范畴。一个 \textbf{$\OO$-幺半函子}是定义在 $\OO^{\otimes}$ 之上、且保持所有余笛卡尔箭头的函子 $\cC^{\otimes}\to\cC^{\prime\otimes}$。令
\[
\Fun^{\otimes}_{\OO}(\cC,\cC')\subseteq\Fun_{\OO^{\otimes}}(\cC^{\otimes},\cC^{\prime\otimes})
\]
为 $\OO$-幺半函子构成的全子 $\infty$-范畴。

令
\[
\Mon_{\OO}(\Cat_{\infty})\subseteq(\Cat_{\infty})_{/\OO^{\otimes}}
\]
为由小 $\OO$-幺半 $\infty$-范畴及其间的 $\OO$-幺半函子构成的 \textbf{1-全}子 $\infty$-范畴。
\end{definition}

\begin{exercise}\label{L21-2.19}
证明：当 $\OO^{\otimes}=\mathrm{Comm}^{\otimes}$ 或 $\mathrm{Assoc}^{\otimes}$ 时，上述定义经由拉直恢复（对称）幺半 $\infty$-范畴与幺半函子的概念。
\end{exercise}

\begin{exercise}\label{L21-2.20}
$\Mon_{\OO}(\Cat_{\infty})$ 中的始对象与终对象分别是什么？
\end{exercise}

\section{代数}

\nitem\label{L21-3.1} 通过余笛卡尔纤维化定义 $\OO$-幺半 $\infty$-范畴的一个好处是，其中代数的定义极其简单。

\begin{definition}\label{L21-3.2}
设 $\OO^{\otimes}$ 与 $\cC^{\otimes}$ 是 $\infty$-操作子。它们之间的一个 \textbf{$\infty$-操作子映射}是定义在 $\mathrm{Fin}_{*}$ 之上、且保持惰性态射的函子 $p\colon\OO^{\otimes}\to\cC^{\otimes}$。我们也称这样的函子为 $\cC$ 中的一个 \textbf{$\OO$-代数}。令
\[
\mathrm{Op}_{\infty}\subseteq(\Cat_{\infty})_{/\mathrm{Fin}_{*}}
\]
为由小 $\infty$-操作子及其间的 $\infty$-操作子映射构成的 \textbf{1-全}子 $\infty$-范畴。

令
\[
\Alg_{\OO}(\cC)\subseteq\Fun_{\mathrm{Fin}_{*}}(\OO^{\otimes},\cC^{\otimes})
\]
为 $\infty$-操作子映射 $\OO^{\otimes}\to\cC^{\otimes}$ 构成的全子 $\infty$-范畴。
\end{definition}

\begin{exercise}\label{L21-3.3}
$\mathrm{Op}_{\infty}$ 中的始对象与终对象分别是什么？
\end{exercise}

\nitem\label{L21-3.4} 要在范畴中定义结合代数，我们只需给它配备幺半结构，而不必是对称幺半结构。更一般地，我们有定义~\ref{L21-3.2} 的下述相对版本。

\begin{definition}\label{L21-3.5}
设
\[
\begin{tikzcd}[column sep=3.4em,row sep=3em]
\OO^{\prime\otimes}\arrow[dr,"q"] & & \cC^{\otimes}\arrow[dl] \\
& \OO^{\otimes} &
\end{tikzcd}
\]
是 $\infty$-操作子之间的映射。$\cC$ 中的一个 \textbf{$\OO'$-代数}（相对于 $\OO$）是
\begin{equation}\label{L21-3.1eq}
\Alg_{\OO'/\OO}(\cC)\defeq\Alg_{\OO'}(\cC)\times_{\Alg_{\OO'}(\OO)}\{q\}
\end{equation}
中的一个对象。当 $q$ 是 $\OO^{\otimes}$ 上的恒等映射时，我们还记
\[
\Alg_{/\OO}(\cC)\defeq\Alg_{\OO/\OO}(\cC).
\]
\end{definition}

\begin{warning}\label{L21-3.6}
$\infty$-范畴 $\Alg_{/\OO}(\cC)$ 与 $\Alg_{\OO}(\cC)$ 是不同的。
\end{warning}

\begin{remark}\label{L21-3.7}
若 $\cC^{\otimes}\to\OO^{\otimes}$ 实现为范畴纤维化，则纤维积~\eqref{L21-3.1eq} 可以按朴素纤维积计算。
\end{remark}

\begin{exercise}\label{L21-3.8}
证明：
\[
(\cC\times_{\OO}\OO')^{\otimes}\defeq\cC^{\otimes}\times_{\OO^{\otimes}}\OO^{\prime\otimes}
\]
是 $\infty$-操作子，且
\[
\Alg_{\OO'/\OO}(\cC)\simeq\Alg_{/\OO'}(\cC\times_{\OO}\OO').
\]
\end{exercise}

\begin{exercise}\label{L21-3.9}
证明：沿 $\langle 1\rangle\in\mathrm{Fin}_{*}$ 的限制给出保守函子
\[
\Alg_{\OO'/\OO}(\cC)\to\Fun_{\OO}(\OO',\cC),
\]
它称为\textbf{遗忘函子}。
\end{exercise}

\begin{definition}\label{L21-3.10}
设 $\cC^{\otimes}$ 与 $\cC^{\prime\otimes}$ 是 $\OO$-幺半 $\infty$-范畴。$\Alg_{\cC/\OO}(\cC')$ 中的一个对象称为从 $\cC^{\otimes}$ 到 $\cC^{\prime\otimes}$ 的（右）\textbf{lax $\OO$-幺半函子}。
\end{definition}

\begin{remark}\label{L21-3.11}
证明：当 $\OO^{\otimes}=\mathrm{Comm}^{\otimes}$ 或 $\mathrm{Assoc}^{\otimes}$ 时，上述定义推广了普通范畴之间的 lax（对称）幺半函子的概念。
\end{remark}

\section{单位性}

\nitem\label{L21-4.1} 注意，对任意

\yizhu{原文此处仅有“Note that for any”，句子明显不完整，底稿疑有脱漏；译文照录，未作补全。\textbf{[译者按（无外部依据）]}}

我们说一个 $\infty$-操作子 $\OO^{\otimes}$ 具有单位性，若对任意颜色 $X\in\OO$，都存在一个本质唯一的零元运算以 $X$ 为目标。

\appendixsection{幺半包络}

\nitem\label{L21-A.1} 按定义，任何对称幺半 $\infty$-范畴都是 $\infty$-操作子。反之则不然，因为 $\infty$-操作子 $\cC^{\otimes}\to\mathrm{Fin}_{*}$ 可能没有足够多的余笛卡尔箭头。

\begin{exercise}\label{L21-A.2}
设 $\cC^{\otimes}\to\mathrm{Fin}_{*}$ 是任意对称幺半 $\infty$-范畴，$\cC'\subseteq\cC$ 是全子 $\infty$-范畴。令 $\cC^{\prime\otimes}\subset\cC^{\otimes}$ 为由满足 $X_i\in\cC'$ 的对象 $(X_1,\cdots,X_n)$ 构成的全子 $\infty$-范畴。证明 $\cC^{\prime\otimes}$ 是 $\infty$-操作子。
\end{exercise}

\begin{exercise}\label{L21-A.3}
反之，证明任何 $\infty$-操作子 $\cC^{\otimes}$ 都可以实现为某个对称幺半 $\infty$-范畴的全子 $\infty$-范畴。
\end{exercise}

\nitem\label{L21-A.4} 事实上，存在一个包含 $\cC^{\otimes}$ 的万有对称幺半 $\infty$-范畴 $\operatorname{Env}(\cC)^{\otimes}$，使得对任意测试用的对称幺半 $\infty$-范畴 $\cD^{\otimes}$，有
\[
\Fun^{\otimes}(\operatorname{Env}(\cC),\cD)\simeq\Alg_{\cC}(\cD).
\]
它称为 $\cC$ 的对称幺半包络。

\nitem\label{L21-A.5} \textbf{建议阅读}\quad HA.2.2.4。
